\begin{tikzpicture}
\begin{semilogxaxis}[
    xlabel={Frecuencia $\omega$ (rad/s)},
    ylabel={Magnitud (dB)},
    grid=both,
    grid style={line width=.1pt, draw=gray!30},
    major grid style={line width=.2pt, draw=gray!60},
    xmin=1e-1, xmax=1e8,
    ymin=-120, ymax=10,
    width=0.9\columnwidth,
    height=5.5cm,
    legend style={
        at={(0.5, 1.05)}, % x=0.5 (centro horizontal), y=1.05 (por encima del límite superior)
        anchor = south,
        scale = 0.9,
        font=\footnotesize % Reduce el tamaño del texto y los símbolos
    }
]

% Trazado asintótico
\addplot[thick, blue] coordinates {
    (1e-1, 0)
    (3.14, 0)
    (3143, -60)
    (3.14e6, 0)
    (1e8, 0)
};
\addlegendentry{Teórico (Asintótico)}

% Marcador de wp1
\draw[dashed, red, thin] (axis cs:3.14, 0) -- (axis cs:3.14, -60);
\addplot[mark=*, mark size=1.5pt, red, forget plot] coordinates {(3.14, -60)};
\node[anchor=north, font=\scriptsize, text=black] at (axis cs:3.14, -60) {$\omega_{p1} = 3.14$};

% Marcador de w0
\addplot[mark=*, mark size=1.5pt, red, forget plot] coordinates {(3143, -60)};
\node[anchor=south, font=\scriptsize, text=black] at (axis cs:3143, -50) {$\omega_{0} = 3143$};

% Marcador de wp2
\draw[dashed, red, thin] (axis cs:3.14e6, 0) -- (axis cs:3.14e6, -60);
\addplot[mark=*, mark size=1.5pt, red, forget plot] coordinates {(3.14e6, -60)};
\node[anchor=north, font=\scriptsize, text=black] at (axis cs:3.14e6, -60) {$\omega_{p2} = 3.14\text{ M}$};

% Datos de LTSpice
\addplot[thick, dashed, orange] table [
    skip first n=1,
    x expr=\thisrowno{0} * 2 * pi,
    y index=1
] {sim/bode_2.txt};
\addlegendentry{Simulación LTSpice}

\end{semilogxaxis}
\end{tikzpicture}
